\begin{textsample}{POS Dim 2 – ro – Score 35.00 – ro/ro\_em\_109.txt}  \label{ex:f2_pos_182}
4.2 Educação de Jovens e Adultos A Educação de Jovens e Adultos – EJA é uma modalidade da Educação Básica que oferece oportunidade de estudos nas etapas do Ensino Fundamental e Médio àqueles que não tiveram acesso ou continuidade na idade própria, nos termos da Lei de \textbf{Diretrizes} e Bases da Educação – LDB/1996, direitos contemplados nos \textbf{artigos} 4º, 5º, 37º, 38º e 87º.
O texto da LDB traz, no § 3º do \textbf{artigo} 37, uma proposição de novos formatos de \textbf{oferta} da EJA : “ A Educação de jovens e adultos deverá articular -se, \textbf{preferencialmente}, com a educação profissional, na forma do regulamento ”.
Diante disso, constata -se que o público da EJA pode ser constituído por \textbf{trabalhadores} que não tiveram a oportunidade de \textbf{cursar} a Educação Básica na infância e adolescência e poderão já estar integrados, em empregos formais ou informais.
Assim, ao retornarem à escola, demandam a especificidade de uma proposta curricular de natureza formativa que \textbf{atenda} tanto a escolarização básica como, também, a formação para o mundo do trabalho.
Além disso, a \textbf{oferta} da EJA atende as \textbf{peculiaridades} de tempo, espaço e rotatividade da população em situação de privação de liberdade, por meio da Educação nos Sistemas Prisional e Socioeducativo, levando em consideração a flexibilidade \textbf{prevista} no \textbf{art}. 23 da Lei nº 9.394/96 ( LDB ) e também estará associada à \textbf{qualificação} profissional, articulando -as, de maneira \textbf{intersetorial}, a \textbf{políticas} e programas \textbf{destinados} a jovens e adultos, e ainda às ações complementares de cultura, esporte, inclusão digital, educação profissional e fomento à leitura.
Nesse contexto, para garantir o direito de jovens e adultos à Educação Básica, o \textbf{currículo} para essa modalidade deve considerar as suas singularidades na elaboração de metodologias e práticas pedagógicas, conforme \textbf{Parecer} CNE/CEB nº 11, de 10 de maio de 2000, a Resolução CNE/CEB nº 1, de 5 de \textbf{julho} de 2000, que estabeleceu as \textbf{Diretrizes} Curriculares Nacionais para a Educação de Jovens e Adultos ( EJA ) ; Resolução CNE/CEB nº 3, de 15 de junho de 2010 ; \textbf{Parecer} CNE/CEB nº 6, de 10 de dezembro de 2020, \textbf{Parecer} nº 1/2021, de 18 de março de 2021 ; e Resolução nº 1/2021, de 25 de maio de 2021, que \textbf{institui} as \textbf{Diretrizes} Operacionais para a EJA alinhadas à Base Nacional Comum Curricular.
Nesse viés, nota -se a necessidade de alinhamento dos \textbf{currículos} e propostas pedagógicas às novas \textbf{legislações} e normas para atendimento a essa modalidade, como dispõe a Resolução nº 1/2021, em seu Art. 13 : > Os \textbf{currículos} dos \textbf{cursos} da EJA, independente de segmento e forma de \textbf{oferta}, deverão garantir, na sua parte relativa à formação geral básica, os direitos e objetivos de aprendizagem, expressos em competências e habilidades nos termos da Política Nacional de Alfabetização ( PNA ) e da BNCC, tendo como ênfase o desenvolvimento dos componentes essenciais para o ensino da leitura e da escrita, assim como das competências gerais e as competências/habilidades relacionadas à Língua Portuguesa, Matemática e Inclusão Digital.
Percebe -se a consonância das \textbf{legislações} \textbf{vigentes} ao Plano Nacional de Educação, aprovado pela Lei nº 13.005, de 25 de junho de 2014, que estabeleceu a Meta 10, definindo que as matrículas de EJA sejam, no mínimo, 25 \% ( vinte e cinco por cento ), nos Ensinos Fundamental e Médio, \textbf{ofertadas} de forma integrada à Educação Profissional.
Com isso, fortalece -se a compreensão de que a modalidade da EJA tem como natureza de \textbf{oferta} o vínculo com a formação profissional e a inserção dos estudantes que a frequentam no mundo do trabalho.
Nota -se pelo último Relatório de Monitoramento da Meta 10 do Plano Nacional de Educação - PNE, correspondente à Meta 12 do Plano Estadual de Educação, que o cumprimento do indicador ainda está aquém do pretendido, apresentando diminuição na \textbf{oferta} da EJA integrada à Educação Profissional em todas as esferas.
Como vemos na trajetória de 2019, cuja previsão era \textbf{atender} 10,39 \% e \textbf{atendeu} apenas 0,2 \% ; e em 2020, com previsão de atingir 12,32 \%, no entanto \textbf{atendeu} 0,1 \%.
Tal fato pode estar relacionado à ausência de dados específicos quanto à modalidade EJA, em relação à \textbf{oferta} de Educação Profissional no formato concomitante e \textbf{cursos} de Formação Inicial e Continuada - FIC.
Com intuito de favorecer o cumprimento das Metas do Plano, ganha destaque a parceria da SEDUC com o Instituto Estadual de Desenvolvimento da Educação Profissional – IDEP, buscando \textbf{atender} as estratégias estabelecidas na referida meta, a fim de avançar na \textbf{oferta} de Educação Profissional para a comunidade em geral, incluindo as formas de EPT articuladas à EJA.
Considerando as necessárias mudanças para atendimento a essa modalidade, a Resolução nº 1/2021 dispõe em seu Art. 2º : > Art. 2º.
Com o objetivo de possibilitar o acesso, a permanência e a continuidade de todas as pessoas que não iniciaram ou interromperam o seu processo educativo escolar, a \textbf{oferta} da modalidade da EJA poderá se dar das seguintes formas : > > I.
Educação de Jovens e Adultos presencial ; > II.
Educação de Jovens e Adultos na modalidade Educação a Distância ( EJA/EaD ) ; > III.
Educação de Jovens e Adultos articulada à Educação Profissional, em \textbf{cursos} de Formação Inicial e Continuada ( FIC ) ou de Formação \textbf{Técnica} de Nível Médio ; > IV.
Educação de Jovens e Adultos com ênfase na Educação e Aprendizagem ao Longo da Vida.
Em consonância às novas \textbf{Diretrizes}, em seu Art. 3, apresenta -se a EJA constituída por meio de regime semestral ou modular, em segmentos e etapas, com a possibilidade de flexibilização do tempo para cumprimento da \textbf{carga} \textbf{horária} exigida, sendo que para cada segmento, há uma correspondência nas etapas da Educação Básica e \textbf{carga} \textbf{horária} específica, destacando -se nesse contexto, a etapa do Ensino Médio, conforme inciso III : > III – para o Ensino médio, que tem como objetivo uma formação geral básica e profissional mais consolidada, seja com a \textbf{oferta} integrada com uma \textbf{qualificação} profissional ou mesmo com um \textbf{curso} \textbf{técnico} de nível médio, \textbf{carga} \textbf{horária} total mínima será de 1.200 ( mil e duzentas ) horas.
Portanto, a etapa do Ensino Médio \textbf{atenderá} pessoas que ainda não concluíram esse segmento, devendo este ser planejado e organizado dentro das novas \textbf{diretrizes} dessa etapa da Educação Básica, contendo uma formação geral básica e a \textbf{oferta} de \textbf{itinerários} tanto propedêuticos quanto técnico-profissionalizantes. \textbf{Preferencialmente}, deve -se buscar uma formação geral profissional mais consolidada, seja com a \textbf{oferta} integrada com uma \textbf{qualificação} profissional ou mesmo com um \textbf{curso} \textbf{técnico} de nível médio.
A formação geral profissional deverá orientar -se pelas demandas cognitivas da área.
Ou seja, um \textbf{curso} de 3º segmento da EJA articulado a uma \textbf{qualificação} profissional devendo aprofundar as competências específicas.
Em \textbf{observância} às \textbf{Diretrizes} Nacionais, o Estado \textbf{normatizou} a \textbf{oferta} da Educação de Jovens e Adultos através das seguintes \textbf{legislações} : Resolução nº 827/10-CEE/RO, que regulamenta a \textbf{oferta} da Educação de Jovens e Adultos no Sistema de Ensino de Rondônia ; \textbf{Portaria} n. 1361/11 GAB/SEDUC, dispõe sobre normas para o funcionamento do \textbf{Curso} Semestral do Ensino Fundamental e Médio – EJA no Estado de Rondônia, \textbf{Portaria} nº 520/2017/SEDUC-NEJA, que estabelece normas para operacionalização da \textbf{oferta} sistemática do \textbf{Curso} Modular na modalidade de Educação de Jovens e Adultos – EJA, \textbf{Portaria} nº 1702/2016-GAB/SEDUC, que implanta o Projeto de Atendimento Diferenciado aos Estudantes no Período Noturno, nas escolas da rede pública estadual de ensino ; e \textbf{Portaria} nº 1165/13-GAB/SEDUC, que dispõe de normas para os \textbf{Exames} Gerais – PROVÃO e de Circulação de Estudos.
Dessa forma, o Referencial Curricular para o Ensino Médio, na modalidade de Educação de Jovens e Adultos, assegura a necessidade de elaboração de matrizes curriculares específicas para esta modalidade, com organização curricular e metodológica diferenciada para os jovens e adultos, considerando as particularidades geracionais, \textbf{preferencialmente} integrada com a formação \textbf{técnica} e profissional, podendo ampliar seus tempos de organização escolar, com menor \textbf{carga} \textbf{horária} diária e anual, garantida a \textbf{carga} \textbf{horária} mínima da parte comum de 1.200 ( um mil e duzentas ) horas e observadas as \textbf{diretrizes} específicas.

% matched lemmas: art, artigo, atender, carga, currículo, cursar, curso, destinar, diretriz, exame, horário, instituir, intersetorial, itinerário, julho, legislação, normatizar, observância, oferta, ofertar, parecer, peculiaridade, política, portaria, preferencialmente, prever, qualificação, trabalhador, técnico, vigente
\end{textsample}
